% Canonical CV source: edit in Overleaf and publish via Integrations > GitHub.
\documentclass[10pt]{article}
\usepackage{geometry}
\geometry{margin=1in}
\usepackage{enumitem}
\setlist[itemize]{topsep=3pt,itemsep=3pt,parsep=0pt,partopsep=0pt}
\usepackage{titlesec}
\usepackage{hyperref}
\usepackage{fancyhdr}
\usepackage{multicol}
\usepackage{parskip}
\usepackage{mathpazo}
\usepackage{csquotes}
\usepackage[normalem]{ulem}
\usepackage{fontawesome5}

% Custom section formatting
\titleformat{\section}{\large\bfseries\uppercase}{\thesection.}{1em}{}
\titleformat{\subsection}[runin]{\bfseries}{}{0pt}{}

% Header/Footer
\pagestyle{fancy}
\fancyhf{}
\fancyfoot[C]{\thepage}
\renewcommand{\headrulewidth}{0pt}

% Custom commands
\newcommand{\sechead}[1]{%
  \vspace{6pt}%
  {\normalsize\bfseries\uppercase{#1}}\\[-8pt]%
  \noindent\rule{\textwidth}{0.4pt}%
  \vspace{2pt}%
}
\newcommand{\subentry}[2]{%
  {\small\textbf{#1}} \hfill {\small\textit{#2}}\\
}
\newcommand{\subsubentry}[2]{%
  {\small#1} \hfill {\small\textit{#2}}\\
}
% Add a consistent gap after the Fields and Thesis lines.
% Change the [9pt] below to tune the space between Education groups.
\newcommand{\subsubentrylast}[2]{%
  {\small#1} \hfill {\small\textit{#2}}\\[9pt]
}
\newcommand{\pubitem}[2]{\item \textit{#1}. #2}
\newcommand{\plainitem}[1]{\item #1}

\begin{document}

% Name and Contact Info
\begin{center}
    {\LARGE \textbf{Daihong Min} \\
    \vspace{0.1em}
    \small\noindent \textit{Updated \today}}\\[5pt]
    \mbox{\small
    \href{https://www.daihongmin.com/}{\faGlobe\ Website} \quad
    \href{https://scholar.google.com/citations?user=9xPqwXAAAAAJ&hl=en}{\faGoogle\ Google Scholar} \quad
    \href{https://www.linkedin.com/in/daihongmin/}{\faLinkedin\ LinkedIn} \quad
    \href{mailto:dhmin@ucdavis.edu}{\faEnvelope\ dhmin@ucdavis.edu} \quad
    \faPhone\ (530) 761-6182}\\[5pt]
    Department of Economics, University of California Davis\\
    One Shields Avenue, Davis, CA 95616
\end{center}

% Sections
\sechead{{\large E}DUCATION}
\subentry{University of California, Davis}{}
\subsubentry{Ph.D. in Economics}{2021--present}
\subsubentrylast{\textit{Fields: International Macroeconomics, International Finance, Political Economics}}{}
\subentry{Korea University}{}
\subsubentry{Ph.D. Program in Economics \textit{(prior to UC Davis)}}{2019--2021}
\subsubentry{M.A. in Economics}{2017--2019}
\subsubentrylast{\textit{Thesis}: ``Effect of Monetary Integration on Purchasing Power Parity in Euro Area"}{}
\subsubentry{B.A. in Economics, \textit{Summa Cum Laude}}{2012--2017}

\sechead{Job Market Paper}
\textbf{Financial Sanctions and Dollar Dominance in Trade Invoicing}\par
\vspace{3pt}
{\small
\textbf{Abstract.} Does the use of financial sanctions weaken dollar dominance in trade invoicing? Using a panel of FX-neutral import invoicing shares for 88 countries over 2012--2023, I examine how responses to U.S. sanctions against Chinese entities vary with baseline import exposure to China. For a one-standard-deviation increase in log sanctions intensity, a 10 percentage points (pp) higher China import share is associated with a contemporaneous Chinese renminbi (RMB) response that is 0.19 pp higher and a U.S. dollar (USD) response that is 0.84 pp lower. The RMB exposure gradient is flatter among destinations whose imports from China have low working-capital intensity. I develop a model in which exporters balance the dollar's normal-time financing and hedging benefits against sanctions-related disruptions to usable credit. The model identifies conditions under which greater working-capital dependence lowers the threshold for switching to RMB invoicing. In the reference quantitative specification, removing working-capital heterogeneity prevents the model from reproducing the low-WCI attenuation, even after re-fitting the remaining parameters. A separate retrospective comparison finds modest improvements in subsequent invoicing paths when working-capital heterogeneity is included. These findings point to the reliability of working-capital finance as a channel through which sanctions can reshape currency choice along trade links to China.\par
}

\sechead{Academic Publications}
\begin{itemize}[leftmargin=*]
    \pubitem{“Purchasing Power Parity vs. Uncovered Interest Rate Parity for NAFTA Countries: The value of Incorporating Time-Varying Parameter Model.”}{\uline{Economic Modelling}, 90 (2020), (with Jong Cheol Yoon and Sang Young Jei)}
    \pubitem{“Empirical Test of Purchasing Power Parity Using a Time-Varying Cointegration Model for China and the UK.”}{\uline{Physica A: Statistical Mechanics and its Applications}, 521 (2019), (with Jong Cheol Yoon and Sang Young Jei)}
    \pubitem{“Impact of Price Disclosure on Gas Market Pricing: Applying Time-varying Cointegration Models.”}{\uline{Journal of Industrial Economics and Business}, 32(5) (2019), (with Jong Cheol Yoon and Sang Young Jei)}
    \pubitem{“Impacts of China's Trade on Sectoral Employment Rates.”}{\uline{Journal of Industrial Economics and Business},  32(3) (2019), (with Jong Chan Lee and Sang Young Jei)}
    \pubitem{“Analysis of the Estimation of Demand Elasticity for Non-Alcoholic Beverage.”}{\uline{Journal of The Korean Data Analysis Society},  21(1) (2019), (with Won Nyon Kim and Sang Young Jei)}
    \pubitem{“China's International Trade and Employment by Sector: Private versus Public.”}{\uline{Journal of The Korean Data Analysis Society},  20(3) (2018), (with Jong Chan Lee and Sang Young Jei)}
\end{itemize}

\newpage
\sechead{Working Papers}
\begin{itemize}[leftmargin=*]
    \plainitem{``Promotional Slack and the Strategic Pass-Through of Local Sugar-Sweetened Beverage Taxes'' (with Yijoong Won and \href{https://kiesel.ucdavis.edu/}{Kristin Kiesel})}
    \plainitem{``Health Monitoring and Household Soda Demand: Evidence from Retail Diabetes-Device Purchases'' (with Yifan Xie)}
\end{itemize}

\sechead{Work in Progress}
\begin{itemize}[leftmargin=*]
    \plainitem{``Tariff News and Inflation Expectations" (with \href{https://sites.google.com/ucdavis.edu/hyunseopark/}{Hyunseo Park} and Yungu Cho)}
    \plainitem{``The Impact of Economic and Monetary Union on Adjustment Speed" (with Sang Young Jei)}
\end{itemize}

\sechead{Grants, Scholarships, and Awards}
\begin{itemize}[leftmargin=*]
    \plainitem{Department Travel Awards, UC Davis}
    \plainitem{Brain Korea 21 Scholarship, Korea University}
    \plainitem{Research Assistant Scholarships, Korea University}
    \plainitem{Teaching Fellow Scholarships, Korea University}
    \plainitem{Best Thesis Awards, Korea University}
    \plainitem{KIMKEEWHA’s Research Scholarships, Korea University}
\end{itemize}

%\sechead{Conference \& Seminar Presentations}
%\begin{itemize}[leftmargin=*]
%    \plainitem{2025: UC Davis}
%\end{itemize}

\sechead{Policy Papers}
\begin{itemize}[leftmargin=*]
    \pubitem{“Analysis of the Effects of Infrastructure Investment on Economic Growth”}{\uline{National Assembly Budget Office}, (2020), (with Yonghun Jung and Lee, Seong-Hoon)} \href{https://www.nabo.go.kr/Sub/01Report/05_Board.jsp?funcSUB=view&bid=19&arg_cid1=0&arg_cid2=0&arg_class_id=0&currentPage=0&pageSize=10&currentPageSUB=0&pageSizeSUB=10&key_typeSUB=subject&keySUB=%EC%9D%B8%ED%94%84%EB%9D%BC&search_start_dateSUB=&search_end_dateSUB=&department=0&department_sub=0&etc_cate1=&etc_cate2=&sortBy=reg_date&ascOrDesc=desc&search_key1=&etc_1=0&etc_2=0&tag_key=&arg_id=7402&item_id=7402&etc_1=0&etc_2=0&name2=0}{(Korean)}
    \pubitem{“The Economic Effects of Gyeongbu-Expressway”}{\uline{Korea Expressway Corporation Research Institute}, (2020), (with Yonghun Jung and Lee, Seong-Hoon)}(Korean)
\end{itemize}

\newpage
\sechead{Teaching Experience}
\subentry{Teaching Assistant, UC Davis}{2021--present}
\vspace{-1em}
\begin{itemize}[leftmargin=*]
    \raggedright
    \plainitem{The Principle of Microeconomics (ECN 1A), The Principle of Macroeconomics (ECN 1B), Intermediate Micro Theory (ECN 100), Intermediate Macro Theory (ECN 101), Economics of International Immigration (ECN 117), Public Finance (ECN 131), Financial Economics (ECN 134), Money, Banks, \& Financial Institutions (ECN 135)}
\end{itemize}
\vspace{1em}
\subentry{Teaching Assistant, Korea University}{2017--2021}
\vspace{-1em}
\begin{itemize}[leftmargin=*]
    \raggedright
    \plainitem{Mathematics for Economics \& Business (ECOP180), Statistics for Economists (ECOP210), Theory and Practice of Public Finance (ECOP303), Econometrics (ECOP305), Taxation and National Economy (ECOP308), Introduction to Financial Times Series Data Analysis (ECOP406), Introduction to Economic Fluctuations (ECOP409)}
\end{itemize}

\sechead{References}
\begin{tabular}{@{}p{0.5\textwidth} p{0.5\textwidth}}
    \textbf{\href{https://faculty.econ.ucdavis.edu/faculty/bergin/}{Paul Bergin}} (Chair) &
    \textbf{\href{https://katherynruss.weebly.com/}{Katheryn Russ}} \\
    Professor & Professor \\
    Department of Economics & Department of Economics \\
    University of California, Davis & University of California, Davis \\
    \href{mailto:prbergin@ucdavis.edu}{\faEnvelope\ prbergin@ucdavis.edu} &
    \href{mailto:knruss@ucdavis.edu}{\faEnvelope\ knruss@ucdavis.edu} \\
    \\[5pt]
    \textbf{\href{https://www.inasimonovska.com/}{Ina Simonovska}} &
    \textbf{\href{https://sites.google.com/view/emilemarin/home?authuser=0}{Emile A. Marin}}  \\
    Associate Professor & Assistant Professor \\
    Department of Economics & Department of Economics \\
    University of California, Davis & University of California, Davis \\
    \href{mailto:inasimonovska@ucdavis.edu}{\faEnvelope\ inasimonovska@ucdavis.edu} &
    \href{mailto:emarin@ucdavis.edu}{\faEnvelope\ emarin@ucdavis.edu}
\end{tabular}

\sechead{Personal Information}
\textbf{Citizenship:} Republic of Korea \\
\textbf{Language:} Korean (native), English (fluent) \\
\textbf{Programming skills:} Stata, Python, R, Matlab, \LaTeX \\

\vfill
\end{document}
